% Created 2025-09-18 jue 09:27
% Intended LaTeX compiler: pdflatex
\documentclass[a4paper]{article}
\usepackage[utf8]{inputenc}
\usepackage[T1]{fontenc}
\usepackage{graphicx}
\usepackage{longtable}
\usepackage{wrapfig}
\usepackage{rotating}
\usepackage[normalem]{ulem}
\usepackage{amsmath}
\usepackage{amssymb}
\usepackage{capt-of}
\usepackage{hyperref}
\usepackage[spanish]{babel}
\usepackage[usenames,dvipsnames]{color} % Required for custom colors
\renewcommand{\ttdefault}{pcr} % MONOESPACIO CON NEGRIT
\usepackage{lastpage}
\usepackage{listings}
\usepackage{listingsutf8}
\renewcommand{\lstlistingname}{Listado}
\lstset{frame=single,inputencoding=utf8,basicstyle=\scriptsize\ttfamily,showstringspaces=false,numbers=none}
\definecolor{MyDarkGreen}{rgb}{0.0,0.4,0.0} % This is the color used for comments
\lstset{ breaklines=true, postbreak=\mbox{\textcolor{red}{$\hookrightarrow$}\space}, keywordstyle=\bfseries, keywordstyle=[1]\color{Blue}\bfseries,  keywordstyle=[2]\color{Purple}\bfseries,  keywordstyle=[3]\color{Blue}\underbar,   identifierstyle=,   commentstyle=\usefont{T1}{pcr}{m}{sl}\color{MyDarkGreen}\small,   stringstyle=\color{Purple},   showstringspaces=false,   tabsize=2,   morecomment=[l][\color{Blue}]{...} }
\lstset{literate=  {á}{{\'a}}1 {é}{{\'e}}1 {í}{{\'i}}1 {ó}{{\'o}}1 {ú}{{\'u}}1   {Á}{{\'A}}1 {É}{{\'E}}1 {Í}{{\'I}}1 {Ó}{{\'O}}1 {Ú}{{\'U}}1   {à}{{\`a}}1 {è}{{\`e}}1 {ì}{{\`i}}1 {ò}{{\`o}}1 {ù}{{\`u}}1   {À}{{\`A}}1 {È}{{\'E}}1 {Ì}{{\`I}}1 {Ò}{{\`O}}1 {Ù}{{\`U}}1   {ä}{{\"a}}1 {ë}{{\"e}}1 {ï}{{\"i}}1 {ö}{{\"o}}1 {ü}{{\"u}}1   {Ä}{{\"A}}1 {Ë}{{\"E}}1 {Ï}{{\"I}}1 {Ö}{{\"O}}1 {Ü}{{\"U}}1   {â}{{\^a}}1 {ê}{{\^e}}1 {î}{{\^i}}1 {ô}{{\^o}}1 {û}{{\^u}}1   {Â}{{\^A}}1 {Ê}{{\^E}}1 {Î}{{\^I}}1 {Ô}{{\^O}}1 {Û}{{\^U}}1   {œ}{{\oe}}1 {Œ}{{\OE}}1 {æ}{{\ae}}1 {Æ}{{\AE}}1 {ß}{{\ss}}1   {ű}{{\H{u}}}1 {Ű}{{\H{U}}}1 {ő}{{\H{o}}}1 {Ő}{{\H{O}}}1   {ç}{{\c c}}1 {Ç}{{\c C}}1 {ø}{{\o}}1 {å}{{\r a}}1 {Å}{{\r A}}1   {€}{{\euro}}1 {£}{{\pounds}}1 {«}{{\guillemotleft}}1   {»}{{\guillemotright}}1 {ñ}{{\~n}}1 {Ñ}{{\~N}}1 {¿}{{?`}}1 }
\usepackage{caption}
\usepackage{attachfile2}
\usepackage[margin=1.5cm,includeheadfoot,includehead,includefoot]{geometry}
\hypersetup{colorlinks,linkcolor=black}
\usepackage{fancyhdr}
\pagestyle{fancyplain}
\chead{}
\lhead{}
\rhead{}
\cfoot{}
\lfoot{\begin{footnotesize}alvaro.gonzalezsotillo@educa.madrid.org\end{footnotesize}}
\rfoot{\begin{footnotesize}\thepage / \pageref{LastPage}\end{footnotesize}}
\usepackage{svg}
\usepackage{letltxmacro}
\LetLtxMacro{\originalincludegraphics}{\includegraphics}
\renewcommand{\includegraphics}[2][]{\IfFileExists{#2.pdf}{\originalincludegraphics[#1]{#2.pdf}}{\originalincludegraphics[#1]{#2}}}
\LetLtxMacro{\originalincludesvg}{\includesvg}
\renewcommand{\includesvg}[2][]{\IfFileExists{#2.pdf}{\originalincludegraphics[#1]{#2.pdf}}{\originalincludegraphics[#1]{#2.svg.pdf}}}
\usepackage{comment}
\excludecomment{NOTES}
\author{Álvaro González Sotillo}
\date{\today}
\title{ENTORNOS DE DESARROLLO\\\medskip
\large  (CÓDIGO: 0487)}
\hypersetup{
 pdfauthor={Álvaro González Sotillo},
 pdftitle={ENTORNOS DE DESARROLLO},
 pdfkeywords={ 0487},
 pdfsubject={},
 pdfcreator={Emacs 30.1 (Org mode 9.7.11)}, 
 pdflang={Spanish}}
\begin{document}

\maketitle
\setcounter{tocdepth}{1}
\tableofcontents

\captionsetup{font=scriptsize}
\section{Cómo serán las clases}
\label{sec:org0000000}
\begin{itemize}
\item Teoría
\begin{itemize}
\item Basada en apuntes
\item Y un aula virtual
\end{itemize}
\item Ejercicios
\begin{itemize}
\item Se realizan en clase o en casa
\item Se ponen en común al día siguiente
\end{itemize}
\item Práctica
\begin{itemize}
\item Máquinas virtuales
\item Posiblemente, una nube
\end{itemize}
\item Trabajos
\begin{itemize}
\item Expondremos algunos en clase
\end{itemize}
\end{itemize}
\section{Materiales}
\label{sec:org0000003}
\begin{itemize}
\item Memoria USB
\item Correo electrónico
\item Acceso a Internet fuera del aula
\end{itemize}
\section{Entrega de Trabajos}
\label{sec:org0000006}
\begin{itemize}
\item Via \textbf{Moodle}
\begin{itemize}
\item Nuestro curso es \url{https://aulavirtual3.educa.madrid.org/ies.alonsodeavellan.alcala}
\item El curso es accesible incluso sin usuario, inicialmente
\end{itemize}
\item Se utilizará Microsoft Office (, \textbf{DOCX}, \textbf{PPTX})
\begin{itemize}
\item Opcionalmente, \textbf{PDF} o LibreOffice (\textbf{ODT}, \textbf{ODP})
\item También es desable entregas \emph{online} (Canvas, Figlet\ldots{})
\end{itemize}
\item Se tendrá en cuenta
\begin{itemize}
\item La corrección técnica de los trabajos
\item La fecha de entrega
\item Expresión, sintaxis, ortografía
\item La apariencia profesional
\end{itemize}
\end{itemize}
\section{Normas}
\label{sec:org000000c}
\begin{itemize}
\item Retrasos y faltas
\item Uso de los ordenadores
\begin{itemize}
\item No pueden utilizarse para tareas distintas de las encargadas por el profesor
\item Se respetará a otros alumnos
\end{itemize}
\item Móviles
\begin{itemize}
\item No.
\end{itemize}
\end{itemize}
\subsection{Averías de los ordenadores}
\label{sec:org0000009}
\begin{itemize}
\item Los problemas se comunican al profesor en cuanto se detectan
\item Se deben hacer copias de seguridad para no perder los datos de los discos
\begin{itemize}
\item Pen Drive
\item Disco Externo
\item Correos enviados a uno mismo
\item Copias en los ordenadores de otros compañeros
\end{itemize}
\item Norma fundamental:
\end{itemize}
\textbf{Si se pierde porque no hay copia, es que no era importante}
\section{Cómo será la evaluación}
\label{sec:org000000f}
\begin{itemize}
\item Las notas de las evaluaciones (1ª,2ª,3ª) no son realmente importantes
\item Solo interesa la nota de la evaluación final
\item Basado en \emph{Resultados de aprendizaje} (RA)
\begin{itemize}
\item Cada RA supone un porcentaje de la nota final
\item Cada prueba (examen, trabajo) indicará que RA evalúa, en qué porcentaje
\item Se necesita aprobar cada RA para aprobar el módulo
\end{itemize}
\end{itemize}
\section{Actividades}
\label{sec:org0000012}
\begin{itemize}
\item Trabajos
\item Exámenes
\item Examen final evaluación ordinaria
\begin{itemize}
\item Con los RA no superados
\end{itemize}
\item Examen evaluación extraordinaria
\begin{itemize}
\item Incluirá todos los RA. La nota del examen será la nota del módulo.
\end{itemize}
\item Entrega de trabajos
\begin{itemize}
\item Individuales, o por parejas si se comparte ordenador
\item Un trabajo entregado fuera de plazo tiene una nota máxima de \emph{6}
\end{itemize}
\end{itemize}
\section{Contenidos}
\label{sec:org0000015}

Según el Decreto 3/2011, de 13 de enero


Desarrollo de software:
\begin{itemize}
\item Concepto de programa informático. Instrucciones y datos.
\item Ejecución de programas en ordenadores:
\begin{itemize}
\item Datos y programas.
\item Hardware vs. software.
\item Estructura funcional de un ordenador: procesador, memoria.
\item Tipos de software. BIOS. Sistema. Aplicaciones.
\item Código fuente, código objeto y código ejecutable; máquinas virtuales.
\end{itemize}
\item Lenguajes de programación:
\begin{itemize}
\item Tipos de lenguajes de programación.
\item Características de los lenguajes más difundidos.
\end{itemize}
\item Introducción a la ingeniería del software:
\begin{itemize}
\item Proceso software y ciclo de vida del software.
\item Fases del desarrollo de una aplicación: análisis, diseño, codificación, pruebas, documentación, explotación y mantenimiento, entre otras.
\begin{itemize}
\item Modelos de proceso de desarrollo software (cascada, iterativo, evolutivo).
\item Metodologías de desarrollo software. Características. Técnicas. Objetivos. Tipos de metodologías:
\item Herramientas CASE (Computer Aided Software Engineering).
\end{itemize}
\end{itemize}
\item Proceso de obtención de código ejecutable a partir del código fuente; herramientas implicadas (editores, compiladores, enlazadores, etcétera).
\item Errores en el desarrollo de programas.
\item Importancia de la reutilización de código.
\end{itemize}

Instalación y uso de entornos de desarrollo:
\begin{itemize}
\item Funciones de un entorno de desarrollo.
\item Tipos de entornos de desarrollo. Entornos de desarrollo libres y propietarios. Características.
\item Instalación de un entorno de desarrollo.
\item Uso básico de un entorno de desarrollo:
\begin{itemize}
\item Uso de herramientas y asistentes.
\item Creación de proyectos.
\item Incorporación de elementos a proyectos.
\item Edición de programas. Sintaxis y formateo de código.
\item Compilación de programas. Detección de errores.
\item Generación de ejecutables.
\item Ejecución de programas.
\item Paneles y vistas.
\item Importación y exportación de ficheros.
\item Personalización.
\item Acceso a documentación.
\item Instalación y desinstalación de aplicaciones, módulos y plugins adicionales.
\item Configuración de actualizaciones.
\item Automatización de tareas.
\end{itemize}
\end{itemize}

Diseño y realización de pruebas:
\begin{itemize}
\item Pruebas en el proceso de desarrollo de software:
\begin{itemize}
\item Planificación de pruebas a lo largo del ciclo de desarrollo.
\item Tipos de pruebas: funcionales, estructurales, regresión, caja negra.
\item Procedimientos y casos de prueba.
\end{itemize}
\item Pruebas de código:
\begin{itemize}
\item Cubrimiento, valores límite, clases de equivalencia.
\item Pruebas unitarias de clases y funciones.
\item Uso de herramientas integradas en los entornos de desarrollo para realizar pruebas unitarias.
\item Automatización de pruebas unitarias.
\item Pruebas de integración.
\item Diseño y documentación casos de prueba.
\end{itemize}
\item Depuración de programas:
\begin{itemize}
\item Herramientas de depuración integradas en los entornos de desarrollo,
\item Puntos de ruptura y seguimiento en tiempo de ejecución.
\item Examinadores de variables.
\end{itemize}
\end{itemize}
Optimización y documentación:
\begin{itemize}
\item Refactorización:
\begin{itemize}
\item Concepto. Limitaciones.
\item Patrones de refactorización más usuales.
\item Refactorización y pruebas.
\item Herramientas de ayuda a la refactorización.
\end{itemize}
\item Control de versiones:
\begin{itemize}
\item Desarrollos colectivos
\item Herramientas de control de versiones. Utilidad. Características. Estructura (cliente/servidor). Repositorio.
\item Clientes de control de versiones. Descarga de ficheros inicial. Modificación de ficheros. Actualización de ficheros en local. Actualización de ficheros en el repositorio. Diferencias entre versiones. Restauración de versiones anteriores. Resolución de conflictos. Historial de versiones.
\end{itemize}

\item Documentación:
\begin{itemize}
\item Uso de comentarios.
\item Herramientas integradas en el entorno de desarrollo para generar documentación automáticas de clases.
\item Alternativas.
\end{itemize}
\end{itemize}

Introducción al lenguaje unificado de modelado (UML, Unified Modeling Language):
\begin{itemize}
\item Características.
\item Versiones.
\item Diagramas UML.
\item Utilización en metodologías de desarrollo orientado a objetos.
\item Herramientas CASE con soporte UML.
\end{itemize}

Elaboración de diagramas de clases:
\begin{itemize}
\item Notación de los diagramas de clases:
\begin{itemize}
\item Clases. Atributos, métodos y visibilidad.
\item Objetos. Instanciación.
\item Relaciones. Asociación, herencia, composición, agregación, dependencia, navegabilidad.
\item Clases abstractas. Interfaces.
\item Paquetes.
\item Grado de detalle.
\end{itemize}
\item Utilización de herramientas CASE para elaborar diagramas de clases.
\item Módulos integrados en entornos de desarrollo para elaborar diagramas de clases.
\item Creación de código a partir de diagramas de clases.
\item Generación de diagramas de clases a partir de código (ingeniería inversa).
\end{itemize}

Elaboración de diagramas de comportamiento:
\begin{itemize}
\item Tipos. Campo de aplicación.
\item Diagramas de casos de uso. Actores, casos de uso, escenario, asociaciones (rela ción de comunicación entre actores y casos de uso), relaciones entre casos de uso.
\item Diagramas de secuencia. Línea de vida de un objeto/actor, activación, envío de mensajes.
\item Diagramas de colaboración. Objetos/actores, mensajes.
\item Otros diagramas:
\begin{itemize}
\item Diagramas de actividades.
\item Diagramas de estado.
\end{itemize}
\item Utilización de herramientas CASE para elaborar diagramas de comportamiento
\item Módulos integrados en entornos de desarrollo para elaborar diagramas de comportamiento.
\end{itemize}
\section{Criterios de evaluación}
\label{sec:org000002a}

Según Real Decreto 450/2010, de 16 de abril
Modificados en Real Decreto 405/2023, de 29 de mayo
\subsection{RA1 Reconoce los elementos y herramientas que intervienen en el desarrollo de un programa informático, analizando sus características y las fases en las que actúan hasta llegar a su puesta en funcionamiento.}
\label{sec:org0000018}
Criterios de evaluación:
\begin{enumerate}
\item Se ha reconocido la relación de los programas con los componentes del sistema informático: memoria, procesador, periféricos, entre otros.
\item Se han identificado las fases de desarrollo de una aplicación informática.
\item Se han diferenciado los conceptos de código fuente, objeto y ejecutable.
\item Se han reconocido las características de la generación de código intermedio para su ejecución en máquinas virtuales.
\item Se han clasificado los lenguajes de programación, identificando sus características.
\item Se ha evaluado la funcionalidad ofrecida por las herramientas utilizadas en el desarrollo de software.
\item Se han identificado las características y escenarios de uso de las metodologías ágiles de desarrollo de software.
\end{enumerate}
\subsection{RA2 Evalúa entornos integrados de desarrollo analizando sus características para editar código fuente y generar ejecutables.}
\label{sec:org000001b}
Criterios de evaluación:
\begin{enumerate}
\item Se han instalado entornos de desarrollo, propietarios y libres.
\item Se han añadido y eliminado módulos en el entorno de desarrollo.
\item Se ha personalizado y automatizado el entorno de desarrollo.
\item Se ha configurado el sistema de actualización del entorno de desarrollo.
\item Se han generado ejecutables a partir de código fuente de diferentes lenguajes en un mismo entorno de desarrollo.
\item Se han generado ejecutables a partir de un mismo código fuente con varios entornos de desarrollo.
\item Se han identificado las características comunes y específicas de diversos entornos de desarrollo.
\end{enumerate}
\subsection{RA3 Verifica el funcionamiento de programas diseñando y realizando pruebas.}
\label{sec:org000001e}
Criterios de evaluación:
\begin{enumerate}
\item Se han identificado los diferentes tipos de pruebas.
\item Se han definido casos de prueba.
\item Se han identificado las herramientas de depuración y prueba de aplicaciones ofrecidas por el entorno de desarrollo.
\item Se han utilizado herramientas de depuración para definir puntos de ruptura y seguimiento.
\item Se han utilizado las herramientas de depuración para examinar y modificar el comportamiento de un programa en tiempo de ejecución.
\item Se han efectuado pruebas unitarias de clases y funciones.
\item Se han implementado pruebas automáticas.
\item Se han documentado las incidencias detectadas.
\item Se han utilizado dobles de prueba para aislar los componentes durante las pruebas.
\end{enumerate}
\subsection{RA4 Optimiza código empleando las herramientas disponibles en el entorno de desarrollo.}
\label{sec:org0000021}
Criterios de evaluación:
\begin{enumerate}
\item Se han identificado los patrones de refactorización más usuales.
\item Se han elaborado las pruebas asociadas a la refactorización.
\item Se ha revisado el código fuente usando un analizador de código.
\item Se han identificado las posibilidades de configuración de un analizador de código.
\item Se han aplicado patrones de refactorización con las herramientas que proporciona el entorno de desarrollo.
\item Se ha realizado el control de versiones integrado en el entorno de desarrollo.
\item Se han utilizado herramientas del entorno de desarrollo para documentar las clases.
\item Se han utilizado repositorios remotos para el desarrollo de código colaborativo.
\item Se han utilizado herramientas para la integración continua del código.
\end{enumerate}
\subsection{RA5 Genera diagramas de clases valorando su importancia en el desarrollo de aplicaciones y empleando herramientas específicas.}
\label{sec:org0000024}
Criterios de evaluación:
\begin{enumerate}
\item Se han identificado los conceptos básicos de la programación orientada a objetos.
\item Se han utilizado herramientas para la elaboración de diagramas de clases.
\item Se ha interpretado el significado de diagramas de clases.
\item Se han trazado diagramas de clases a partir de las especificaciones de las mismas.
\item Se ha generado código a partir de un diagrama de clases.
\item Se ha generado un diagrama de clases mediante ingeniería inversa.
\end{enumerate}
\subsection{RA6 Genera diagramas de comportamiento valorando su importancia en el desarrollo de aplicaciones y empleando herramientas específicas.}
\label{sec:org0000027}
Criterios de evaluación:
\begin{enumerate}
\item Se han identificado los distintos tipos de diagramas de comportamiento.
\item Se ha reconocido el significado de los diagramas de casos de uso.
\item Se han interpretado diagramas de interacción.
\item Se han elaborado diagramas de interacción sencillos.
\item Se ha interpretado el significado de diagramas de actividades.
\item Se han elaborado diagramas de actividades sencillos.
\item Se han interpretado diagramas de estados.
\item Se han planteado diagramas de estados sencillos.
\end{enumerate}
\section{Distribución de RA en unidades de trabajo}
\label{sec:org000002d}

\href{ed-0-contenidos-criterios-evaluacion.xlsx}{Enlace a hoja excel}
\section{Referencias}
\label{sec:org0000030}
\phantomsection
\label{}
\begin{itemize}
\item Formatos:
\begin{itemize}
\item \href{./ed-0-contenidos-criterios-evaluacion.reveal.html}{Transparencias}
\item \href{./ed-0-contenidos-criterios-evaluacion.pdf}{PDF}
\item \href{./ed-0-contenidos-criterios-evaluacion.wp.html}{Página web}
\item \href{./ed-0-contenidos-criterios-evaluacion.epub}{EPUB}
\end{itemize}
\item Creado con:
\begin{itemize}
\item \href{https://www.gnu.org/s/emacs/}{Emacs}
\item \href{https://gitlab.com/oer/org-re-reveal}{org-re-reveal}
\item \href{https://www.latex-project.org/}{Latex}
\end{itemize}
\item Alojado en \href{https://alvarogonzalezsotillo.github.io/apuntes-clase}{Github}
\end{itemize}
\end{document}
